% Cognitive Computation manuscript starter
% Based on the Springer Nature LaTeX authoring template, version 3.1
% (December 2024), adapted to the journal's current submission guidelines.
%
% Compile with:
%   latexmk -pdf cognitive_computation_manuscript.tex
%
% The referee option enables double spacing for review. Remove it only if the
% editorial office requests a different presentation.
\documentclass[referee,pdflatex,sn-vancouver-num]{sn-jnl}

\usepackage{graphicx}
\usepackage{amsmath,amssymb,amsfonts}
\usepackage{booktabs}
\usepackage{multirow}

\raggedbottom

\begin{document}

\title[Short running title]{Full Informative Title of the GENEActiv Study}

% Mark the corresponding author with an asterisk.
% Add further authors by repeating \author and linking affiliation numbers.
\author*[1]{\fnm{Given Name} \sur{Family Name}}\email{corresponding.author@example.org}
\author[1]{\fnm{Given Name} \sur{Family Name}}

\affil*[1]{\orgdiv{Department}, \orgname{Institution},
  \orgaddress{\street{Street address}, \city{City}, \postcode{Postal code},
  \country{Country}}}

% The journal requires a structured abstract of 150--250 words using these
% four headings. Do not cite references or use undefined abbreviations here.
\abstract{%
\textbf{Background / introduction:} State the clinical or computational problem,
the knowledge gap, and the objective of the study.

\textbf{Methods:} Summarize the cohort, GENEActiv recording protocol, feature
extraction, outcomes, statistical or machine-learning analysis, validation
strategy, and primary evaluation metrics.

\textbf{Results:} Report the main quantitative findings with uncertainty where
appropriate. Distinguish prespecified analyses from exploratory analyses.

\textbf{Conclusions:} Explain the principal interpretation, relevance to
cognitive computation, and the main limitation or implication.}

% Supply 4--6 indexing terms (the journal page is inconsistent in one place,
% saying 3--6, but its detailed preparation instructions specify 4--6).
\keywords{GENEActiv, actigraphy, wearable sensors, cognitive computation}

\maketitle

\section{Introduction}\label{sec:introduction}

% Keep the introduction concise. Establish the problem, prior work, gap,
% research question or hypotheses, and contribution.
Describe the motivation for analysing GENEActiv measurements and state the
study objective. Cite sources in numerical order using \verb|\cite{key}|; the
selected Vancouver style renders citations in square brackets. For example,
the original GENEA accelerometer validation can be cited as
\cite{esliger2011validation}; replace or retain this source as appropriate.

\section{Methods}\label{sec:methods}

% Methods published in detail elsewhere may be cited rather than repeated.
% Define every abbreviation at first use and use SI units throughout.

\subsection{Study design and participants}\label{subsec:participants}

Describe recruitment, inclusion and exclusion criteria, diagnostic groups,
sample size, and the timing of study visits. State the ethics committee name,
approval identifier, and consent procedure.

\subsection{GENEActiv data acquisition}\label{subsec:acquisition}

Report device model, placement, sampling frequency, recording duration,
calibration, non-wear handling, and any sleep-diary or clinical annotations.

\subsection{Signal processing and feature extraction}\label{subsec:processing}

Document the processing pipeline sufficiently for reproduction: input format,
filtering, epoch definition, sleep or activity segmentation, derived features,
missing-data handling, normalization, and software versions.

\subsection{Clinical and cognitive measures}\label{subsec:clinical}

Define the clinical outcomes, scales, cognitive domains, and group labels used
in the analysis. Explain how measurements were aligned with wearable data.

\subsection{Statistical analysis}\label{subsec:statistics}

Specify primary and secondary analyses, covariates, multiple-comparison
control, effect sizes, confidence intervals, missing-data strategy, and the
significance threshold.

\subsection{Machine-learning analysis}\label{subsec:machine-learning}

If applicable, describe preprocessing within each training fold, participant-
grouped data splitting, model selection, hyperparameter tuning, class-imbalance
handling, leakage controls, and final evaluation metrics.

\section{Results}\label{sec:results}

Present participant flow and cohort characteristics before the primary results.
Report numerical estimates, uncertainty, and exact sample sizes. Refer to every
table and figure in consecutive order.

\begin{table}[t]
\caption{Placeholder for cohort characteristics or primary results. Replace
all example labels and values before submission.}\label{tab:example}
\begin{tabular}{@{}lll@{}}
\toprule
Measure & Group 1 & Group 2 \\
\midrule
Participants, $n$ & -- & -- \\
Age, years & -- & -- \\
Valid recording time, days & -- & -- \\
\bottomrule
\end{tabular}
\end{table}

% Example figure inclusion (uncomment after adding the file beside this .tex):
% \begin{figure}[t]
%   \centering
%   \includegraphics[width=\linewidth]{figure_1.pdf}
%   \caption{A self-contained caption defining all symbols and abbreviations.}
%   \label{fig:example}
% \end{figure}

\section{Discussion}\label{sec:discussion}

Interpret the principal findings in relation to the hypotheses and prior work.
Discuss clinical and computational relevance, robustness analyses, alternative
explanations, generalizability, and limitations without repeating the Results.

\section{Conclusion}\label{sec:conclusion}

State the main conclusion and its implications concisely. Avoid introducing
new results.

\backmatter

\section*{Acknowledgements}

Name contributors who do not meet authorship criteria and acknowledge grants or
other support. Write funding organization names in full.

\section*{Declarations}

% The journal requests the following disclosures as applicable. Replace every
% placeholder; use "Not applicable" only when it is genuinely inapplicable.

\subsection*{Funding}

[State funder names and grant numbers, or ``The authors received no funding for
this work.'']

\subsection*{Competing interests}

[Disclose financial and non-financial interests, or ``The authors have no
relevant financial or non-financial interests to disclose.'']

\subsection*{Ethics approval}

[Name the approving ethics committee or institutional review board, approval
number, and the ethical standard under which the study was conducted.]

\subsection*{Consent to participate}

[State how informed consent was obtained, or explain why it was waived.]

\subsection*{Consent for publication}

[State whether consent for publication is applicable and was obtained.]

\subsection*{Data availability}

[Original research must include this statement. Give repository links and
persistent identifiers, or explain privacy restrictions and the controlled
access procedure.]

\subsection*{Code availability}

[Give the repository URL, archived release DOI, license, and any access
conditions, or explain why code cannot be shared.]

\subsection*{Author contributions}

[List each author's contributions, preferably using the CRediT taxonomy.]

% Keep references.bib and sn-vancouver-num.bst beside this file when uploading.
\bibliography{references}

\end{document}
